% !TeX root = GeoDynamics.tex
\section{Programming Styles}

When discussing software architecture, two common programming styles frequently emerge. Although there is not a universally accepted pair of names for these approaches, they are often described as \emph{control-oriented} and \emph{data-oriented} designs.

\subsection{Control-Oriented (Hierarchical) Design}

In this style, functions determine what happens next by directly invoking other functions within their bodies.

\begin{lstlisting}[language=C++]
void Process()
{
LoadData();

```
if(condition)
    MethodA();
else
    MethodB();
```

}

void MethodA()
{
...
MethodC();
}
\end{lstlisting}

Characteristics include:

\begin{itemize}
\item Deep call trees.
\item Control flow is distributed among functions.
\item Functions decide what operation occurs next.
\item Execution resembles a hierarchy or tree structure.
\end{itemize}

Common names for this style include:

\begin{itemize}
\item Hierarchical Design
\item Top-Down Design
\item Control-Oriented Programming
\item Call-Tree Architecture
\item Procedural Decomposition
\end{itemize}

A potential drawback is that execution flow may become difficult to follow because the program continually branches into other functions.

\subsection*{Call-Tree Representation}

\begin{verbatim}
Function A
|
+-- Function B
|   |
|   +-- Function D
|   |
|   +-- Function E
|
+-- Function C
\end{verbatim}

The primary focus is on \emph{who calls whom}.

\subsection{Data-Oriented (Pipeline) Design}

In this style, each function performs a specific operation and returns data to the caller. The returned data then becomes the input for the next stage.

\begin{lstlisting}[language=C++]
auto data = LoadData();
auto particles = BuildParticles(data);
auto collisions = DetectCollisions(particles);
auto result = ResolveCollisions(collisions);
\end{lstlisting}

Characteristics include:

\begin{itemize}
\item Execution path is explicit.
\item Functions have clear inputs and outputs.
\item Easier testing and verification.
\item Improved modularity and reuse.
\end{itemize}

Common names include:

\begin{itemize}
\item Pipeline Architecture
\item Data-Flow Programming
\item Functional Composition
\item Chain of Transformations
\item Sequential Processing
\end{itemize}

\subsubsection*{Pipeline Representation}

\begin{verbatim}
Data
|
v
Stage A
|
v
Stage B
|
v
Stage C
|
v
Result
\end{verbatim}

The primary focus is on \emph{how data moves through the system}.

\subsection{Application to GPU and HPC Systems}

Many high-performance computing (HPC) and GPU applications naturally adopt a pipeline architecture.

For example, a particle collision detection system may be represented as:

\begin{verbatim}
Particles
|
v
Broad Phase
|
v
Candidate Pairs
|
v
Narrow Phase
|
v
Collision Resolution
|
v
Updated Particles
\end{verbatim}


Each stage consumes data produced by the previous stage and generates data for the next stage.

This architecture is common in:

\begin{itemize}
\item GPU pipelines
\item Vulkan rendering systems
\item Compute shader workflows
\item Scientific computing applications
\item Data processing frameworks
\end{itemize}

\subsection{Summary}

The distinction between these approaches can be summarized as follows:

\begin{center}
\begin{tabular}{|l|l|}
\hline
\textbf{Control-Oriented} & \textbf{Data-Oriented} \\
\hline
Focus on function calls & Focus on data movement \\
Hierarchical call trees & Linear processing stages \\
Control flow distributed & Control flow explicit \\
Often deeper nesting & Often flatter structure \\
Harder to parallelize & Easier to parallelize \\
\hline
\end{tabular}
\end{center}
For large simulation codes, physics engines, and GPU applications, data-oriented pipeline designs are often preferred because they make data dependencies explicit and are generally easier to reason about, profile, parallelize, and maintain.

\subsection{Generated Source Flowcharts}

For the Geo and shadow dynamics work, the sequential stage chain should be
documented with a generated flowchart when the source structure becomes hard to
hold in memory.  The flowchart is not a substitute for narrow functions; it is a
debugging and navigation aid that makes the stage order visible.

The current source-flow generator is:

\begin{verbatim}
python/base/SourceFlowChart.py
\end{verbatim}

The current system-flow generator is:

\begin{verbatim}
python/base/SystemFlowChart.py
\end{verbatim}

It emits DOT, SVG, PDF, and PNG artifacts under
\texttt{docs/Dynamics/generated}.  The DOT file preserves the graph source, the
SVG can carry interactive links, and the PDF is used as the stable LaTeX figure.
The source-flow chart links mostly to code.  The system-flow chart links mostly
to documentation because it explains loops, decision trees, storage boundaries,
and GLSL-shaped algorithm stages.  The generators should be updated when
function names, stage order, decision points, or major source handoffs change.

This preserves the same programming rule used by the code: the primary flow is
explicit, generated documentation follows that flow, and helpers are shown as
supporting calls rather than hidden control-flow owners.

\subsection{Source-Complete Linear Dynamics}

The linear shadow implementation demonstrates the preferred shape for dynamics
that will eventually execute once per source particle in GLSL.  The top-level
function must read like the algorithm:

\begin{verbatim}
initialize source totals
for each source-owned contact:
    initialize one contact
    calculate one contact
calculate source contact weights
calculate planned contact impulses
apply one source velocity write
update source-owned ledgers
\end{verbatim}

Loops belong in the narrowest function that owns the sequence.  Per-contact
functions process exactly one contact and do not hide a second contact-list
loop.  Source-level functions may iterate the completed source contact slots
when their result depends on the whole source contact set.

This shape keeps the main decision order visible, makes stepper debugging
practical, and exposes which functions must eventually map into a per-source
shader invocation.  Movement remains a separate pass because every source must
read the same frame-start positions and velocities during contact response.

\subsection{Interactive Source Documentation}

When the generated flow is included in the project PDF, native TikZ is preferred
for interactive navigation.  Each function node links to a stable labeled
function section in \texttt{SourceDocumentation.tex}.  Graphviz remains useful
for standalone DOT, SVG, PDF, and PNG artifacts, but links embedded in a
Graphviz PDF are not expected to survive inclusion through
\texttt{\textbackslash includegraphics}.

